\section{El Entorno de Ejecución}

Antes de que el sistema ROS~2 pueda ser invocado, el sistema operativo debe saber dónde encontrarlo. Esta sección explica la infraestructura sobre la que todo lo demás descansa.

\subsection{Distribuciones de ROS~2}

ROS~2 sigue un esquema de versiones llamado distribuciones, análogo al de Ubuntu o Debian. Cada distribución tiene un nombre en clave formado por un adjetivo y un animal, ambos iniciando con la misma letra, en orden alfabético. Las principales distribuciones son:

\begin{table}[H]
  \centering
  \caption{Distribuciones de ROS~2 y su soporte.}
  \label{tab:distros}
  \begin{tabular}{lll}
    \toprule
    Distribución & Año & Soporte \\
    \midrule
    Foxy Fitzroy        & 2020 & Expirado \\
    Galactic Geochelone & 2021 & Expirado \\
    Humble Hawksbill    & 2022 & LTS hasta 2027 \\
    Iron Irwini         & 2023 & Expirado \\
    Jazzy Jalisco       & 2024 & LTS hasta 2029 \\
    \bottomrule
  \end{tabular}
\end{table}

Las distribuciones LTS (Long Term Support) tienen soporte garantizado por cinco años y son las recomendadas para proyectos de investigación y producción. \textbf{Humble Hawksbill} es la distribución LTS activa más extendida al momento de este trabajo y la que se utiliza en todos los experimentos descritos.

Cada distribución se instala en su propio directorio dentro de \texttt{/opt/ros/}. Si se tuvieran Humble y Jazzy instaladas simultáneamente, coexistirían en \texttt{/opt/ros/humble/} y \texttt{/opt/ros/jazzy/} sin interferir entre sí. La distribución activa en una terminal es aquella cuyo \texttt{setup.bash} fue sourceado, y la variable \texttt{ROS\_DISTRO} indica cuál es:

\begin{lstlisting}[style=bash, caption={Verificación de la distribución activa.}]
echo $ROS_DISTRO
# humble
\end{lstlisting}

\subsection{Dónde vive ROS~2}

ROS~2 Humble se instala en \texttt{/opt/ros/humble/}. El directorio \texttt{/opt/} es el lugar estándar en Linux para software opcional instalado en el sistema, es decir, software que no viene con el OS pero que se agrega encima. La estructura interna es análoga a la del propio sistema operativo:

\begin{lstlisting}[style=bash, caption={Estructura de instalación de ROS~2 Humble.}]
/opt/ros/humble/
|-- bin/       <- ejecutables: ros2, rviz2, gazebo...
|-- lib/       <- librerias compiladas y paquetes Python
|-- include/   <- headers de C++
|-- share/     <- recursos y launch files del sistema
`-- setup.bash <- el script que se sourcea
\end{lstlisting}

El workspace del usuario replica exactamente esta estructura, pero contiene únicamente los paquetes desarrollados por él:

\begin{lstlisting}[style=bash, caption={Estructura del workspace tras compilar con \texttt{colcon build}.}]
ws3/install/
|-- control_nodes/
|   |-- bin/
|   `-- lib/python3.10/site-packages/
|-- simulador/
`-- setup.bash  <- encadenador: sourcea ROS2 + workspace
\end{lstlisting}

\subsection{Variables de entorno}

Una variable de entorno es un par \texttt{NOMBRE=valor} almacenado en la memoria del proceso actual. No pertenece a ningún programa en particular, sino al entorno del proceso: la tabla de configuración global que el sistema operativo mantiene para cada proceso en ejecución.

\texttt{PATH} es la variable que contiene una lista de directorios separados por \texttt{:} donde el OS busca ejecutables. Si \texttt{/opt/ros/humble/bin} no está en \texttt{PATH}, el comando \texttt{ros2} devuelve \texttt{command not found}. \texttt{PYTHONPATH} cumple la misma función para módulos Python: sin esta variable configurada, la instrucción \texttt{import rclpy} falla con \texttt{ModuleNotFoundError}. \texttt{AMENT\_PREFIX\_PATH} es utilizada internamente por ROS~2 para localizar los paquetes instalados mediante \texttt{ament\_index}. Finalmente, \texttt{ROS\_DISTRO} contiene el nombre de la distribución activa y es leída por algunos paquetes para ajustar su comportamiento. Todas pueden inspeccionarse con:

\begin{lstlisting}[style=bash, caption={Inspección de variables de entorno relevantes.}]
echo $PATH        # rutas donde el OS busca ejecutables
echo $PYTHONPATH  # rutas donde Python busca modulos
echo $ROS_DISTRO  # distribucion de ROS2 activa
env               # muestra todas las variables
\end{lstlisting}

\subsection{Qué es \texttt{source} y por qué es necesario}

Cuando se ejecuta un script de shell de la manera convencional, el OS crea un proceso hijo, ejecuta el script en ese proceso hijo, y al terminar el proceso hijo desaparece llevándose consigo cualquier cambio a las variables de entorno. La terminal del usuario no es afectada.

\begin{lstlisting}[style=bash, caption={Diferencia entre ejecutar y sourcear un script.}]
# Ejecutar normal: modifica un proceso hijo que desaparece
bash setup.bash      # los cambios a PATH se pierden

# Sourcear: ejecuta las lineas directamente en la terminal
source setup.bash    # los cambios a PATH persisten
\end{lstlisting}

El comando \texttt{source} (equivalente al punto \texttt{.} en bash) ejecuta las líneas del script en el proceso actual de la terminal, como si se hubieran escrito directamente. Las modificaciones a \texttt{PATH} y \texttt{PYTHONPATH} quedan en ese proceso y son heredadas por todos los programas que se lancen desde ahí.

\subsection{Qué hace \texttt{setup.bash} internamente}

El archivo \texttt{ws3/install/setup.bash} generado por \texttt{colcon build} actúa como encadenador: primero sourcea ROS~2, luego el workspace del usuario. Su función central es invocar un script Python auxiliar que inspecciona los paquetes instalados en \texttt{ws3/install/} y genera dinámicamente los comandos \texttt{export} necesarios:

\begin{lstlisting}[style=bash, caption={Mecanismo central de \texttt{local\_setup.bash} (simplificado).}]
# El script Python lee ws3/install/ y genera los export
_cmds="$(python3 _local_setup_util_sh.py sh bash)"

# eval ejecuta esas lineas en la terminal actual
eval "${_cmds}"

# El resultado equivale aproximadamente a:
export PATH="/ws3/install/control_nodes/bin:$PATH"
export PYTHONPATH="/ws3/install/control_nodes/lib/\
python3.10/site-packages:$PYTHONPATH"
\end{lstlisting}

Esto es importante: \texttt{source setup.bash} no instala nada ni configura nada de forma permanente. Solo modifica las variables de entorno de la terminal actual. Al cerrar esa terminal, las modificaciones desaparecen. Por eso es necesario sourcear en cada terminal nueva.

\subsection{El shebang y el PATH}

El ejecutable \texttt{ros2} es un archivo de texto. Su primera línea contiene:

\begin{lstlisting}[style=bash, caption={Primera línea del ejecutable \texttt{ros2}.}]
#!/usr/bin/env python3
\end{lstlisting}

Esta línea se llama shebang. Cuando Linux intenta ejecutar un archivo con permisos de ejecución, lee su primera línea: si comienza con \texttt{\#!}, usa lo que sigue como intérprete y le pasa el archivo como argumento. \texttt{/usr/bin/env python3} no apunta a Python directamente; llama a \texttt{env}, que busca \texttt{python3} en el \texttt{PATH} actual. Como ya se sourceó ROS~2, ese \texttt{python3} tiene acceso a todos los módulos de ROS~2 a través de \texttt{PYTHONPATH}.

En consecuencia, escribir \texttt{ros2 launch ...} en la terminal es exactamente equivalente a:

\begin{lstlisting}[style=bash, caption={Equivalencia entre el comando y su expansión real.}]
python3 /opt/ros/humble/bin/ros2 launch \
    simulador demo_pva.launch.py \
    n_rays_used:=5 vmax:=0.8
\end{lstlisting}
